% !TeX root = ../navier-stokes-manuscript.tex
\subsection{What the singularity demands}\label{sub:BodyFiniteTimeResponseEscalation}

The singularity cannot simply appear at \(\tau\).
The velocity has to reach every preceding state, and the intervals available for those changes are collapsing.
Navier--Stokes already tells us how to read the first response.

Along a particle path \(X'(t)=u(X(t),t)\),
\[
  \frac{d}{dt}u(X(t),t)
  =D_tu(X(t),t),
  \qquad
  D_tu:=\partial_tu+(u\cdot\nabla)u.
\]
The momentum equation is therefore an equation for the material acceleration:
\begin{equation}\label{eq:MaterialAccelerationEquationBody}
  D_tu=-\nabla p+\nu\Delta u.
\end{equation}
Applying \(D_t\) again produces the material jerk, but it also differentiates the path and the transport field.
To see the temporal responses that repeat inside this calculation, hold the spatial coordinate fixed and write
\begin{equation}\label{eq:TemporalTowerDefinitionBody}
  V_n:=\partial_t^nu,
  \qquad
  Q_n:=\partial_t^np,
  \qquad
  n\in\Nzero.
\end{equation}
Material acceleration and the first Eulerian row are related by
\[
  D_tu=V_1+(V_0\cdot\nabla)V_0.
\]
They are not the same object.
The rows \(V_0,V_1,V_2,\ldots\) are the fixed-position responses that repeated time differentiation of the equation actually generates.

Now take one of those rows on an interval \([s,t]\).
Absolute continuity in a fixed Banach space \(B\) gives
\[
  V_n(t)-V_n(s)
  =\int_s^tV_{n+1}(\tau)\,d\tau,
\]
and therefore
\begin{equation}\label{eq:BodyResponseEscalation}
  \operatorname*{ess\,sup}_{s<\tau<t}
  \norm{V_{n+1}(\tau)}_B
  \geq
  \frac{\norm{V_n(t)-V_n(s)}_B}{t-s}.
\end{equation}
If the interval shrinks while \(V_n\) still makes a definite change, the next response must grow.
More precisely, if
\[
  t_k-s_k\longrightarrow0,
  \qquad
  \norm{V_n(t_k)-V_n(s_k)}_B\geq c_n>0,
\]
then \(V_{n+1}\) becomes unbounded along those intervals.
The condition has to be paid at every rung: the next derivative is forced only when the current response continues to carry a nontrivial change.

This is the exact content behind the velocity--acceleration--jerk picture.
Velocity changing across a collapsing interval forces a large first response.
If that response must change again across the next collapsing interval, jerk has to lead it.
If the jerk must change, the next derivative has to lead the jerk.
None of the later responses can stagnate while it is still being asked to produce the next state of the cascade.

Now differentiate the equation.
At the first row,
\[
  V_1
  =\nu\Delta V_0
  -(V_0\cdot\nabla)V_0
  -\nabla Q_0.
\]
Differentiate once.
The time derivative can land on either copy of the velocity in the transport term:
\[
  V_2
  =\nu\Delta V_1
  -(V_1\cdot\nabla)V_0
  -(V_0\cdot\nabla)V_1
  -\nabla Q_1.
\]
Differentiate again, and the middle placement appears twice:
\[
  V_3
  =\nu\Delta V_2
  -(V_2\cdot\nabla)V_0
  -2(V_1\cdot\nabla)V_1
  -(V_0\cdot\nabla)V_2
  -\nabla Q_2.
\]
At order \(n\), the binomial coefficients count every way the time derivatives can divide between the transporting velocity and the velocity being transported:
\begin{equation}\label{eq:BodyTemporalVelocityRecurrence}
  V_{n+1}
  +\sum_{a=0}^{n}\binom{n}{a}(V_a\cdot\nabla)V_{n-a}
  +\nabla Q_n
  =\nu\Delta V_n,
  \qquad
  \nabla\cdot V_n=0.
\end{equation}
These differentiated equations are the temporal Tower generated by Navier--Stokes.

Pressure follows the same differentiation.
Taking the divergence of the original equation gives
\begin{equation}\label{eq:PressureNormalization}
  -\Delta p=\partial_i\partial_j(u_i u_j),
  \qquad
  p=R_iR_j(u_i u_j),
\end{equation}
where decay at spatial infinity fixes the additive constant.
At temporal order \(n\),
\begin{equation}\label{eq:BodyTemporalPressureRecurrence}
  Q_n
  =R_iR_j
  \sum_{a=0}^{n}\binom{n}{a}
  (V_a)_i(V_{n-a})_j.
\end{equation}
The velocity rows determine \(Q_n\), and \(\nabla Q_n\) returns immediately in the equation for \(V_{n+1}\).
Pressure is present at every rung.

Each temporal row is spatial as well.
Viscosity contributes \(\Delta V_n\); transport differentiates in space; pressure is reconstructed from spatial singular integrals of the same velocity products.
Following those derivatives introduces the mixed coordinates
\begin{equation}\label{eq:BodyMixedJetDefinition}
  J_{n,\alpha}:=\partial_t^n\partial_x^\alpha u,
  \qquad
  \Pi_{n,\alpha}:=\partial_t^n\partial_x^\alpha p,
  \qquad
  \alpha\in\Nzero^3.
\end{equation}
One spatial derivative already splits the transport again:
\[
  \partial_\ell\bigl((V_a\cdot\nabla)V_{n-a}\bigr)
  =((\partial_\ell V_a)\cdot\nabla)V_{n-a}
  +(V_a\cdot\nabla)\partial_\ell V_{n-a}.
\]
For a general multi-index \(\alpha\), Leibniz's rule gives
\begin{align}
  J_{n+1,\alpha}
  &+
  \sum_{a=0}^{n}
  \sum_{\beta\leq\alpha}
  \binom{n}{a}\binom{\alpha}{\beta}
  (J_{a,\beta}\cdot\nabla)
  J_{n-a,\alpha-\beta}
  +\nabla\Pi_{n,\alpha}
  \notag\\
  &=
  \nu\Delta J_{n,\alpha},
  \qquad
  \nabla\cdot J_{n,\alpha}=0.
  \label{eq:BodyMixedJetRecurrence}
\end{align}
The pressure coordinates satisfy
\begin{equation}\label{eq:BodyMixedPressureRecurrence}
  -\Delta\Pi_{n,\alpha}
  =
  \partial_x^\alpha\partial_i\partial_j
  \sum_{a=0}^{n}\binom{n}{a}
  (V_a)_i(V_{n-a})_j.
\end{equation}

Following the singularity into Navier--Stokes leaves two independent indices.
Moving down changes the temporal response; moving across resolves that response more finely in space; and the equation couples both directions in every row.
The surprise is what viscosity makes visible along the spatial direction.
